\section{Exact weight-enumerator certificates}\label{app:certificates}

This appendix specifies the necessary linear systems used for the
low-weight reductions. The certificate files contain their integer
multipliers and all case assumptions. The checker rebuilds the columns
from physical Pauli cosets.

\subsection{Split transforms and the affine shadow}

The ordinary weight enumerator records only total weight. A split
enumerator also records which coordinate blocks contain the nonidentity
letters. This is needed when a specified stabilizer word or subgroup
occupies one block. The formulas below are the MacWilliams and shadow
identities applied separately to each block.

Partition the physical coordinates into blocks of sizes
\(n_1,\ldots,n_s\). Let \(A^{\rm split}_{\mathbf i}\) count stabilizer words with
split weight \(\mathbf i=(i_1,\ldots,i_s)\). Put \(B=|S|=2^{n-k}\), and
\[
 \mathcal{K}_j^{\ell}(w)=\sum_{u=\max(0,j-\ell+w)}^{\min(j,w)}
          (-1)^u3^{j-u}\binom wu\binom{\ell-w}{j-u}.
\]
Character orthogonality gives the centralizer count
\[
 |\{v\in S^\perp:\operatorname{splitwt}(v)=\mathbf j\}|
 =B^{-1}\mathcal T_{\mathbf j},\quad
 \mathcal T_{\mathbf j}=\sum_{\mathbf i}
       \left(\prod_\ell \mathcal{K}_{j_\ell}^{n_\ell}(i_\ell)\right)A^{\rm split}_{\mathbf i}.
\]
The character identity behind this formula is
\[
 \mathbf 1_{v\in S^\perp}=B^{-1}\sum_{s\in S}(-1)^{\langle v,s\rangle}.
\]
The sum is \(B\) when every word of \(S\) commutes with \(v\),
and zero otherwise. Summing it over words of a fixed split weight
produces the products of Krawtchouk polynomials above. This explains why
the transformed coefficients count actual centralizer words.
On an isotropic stabilizer the parity function is linear, by
\eqref{eq:polar}. The affine set
\[
 \mathcal A=\{v:\langle v,s\rangle=\wt(s)\bmod2
                          \text{ for every }s\in S\}
\]
is a nonempty coset of \(S^\perp\), since every linear functional on
\(S\) extends to the ambient space. Character orthogonality gives its
split weight counts as \(B^{-1}\widehat{\mathcal T}_{\mathbf j}\), where
\[
 \widehat{\mathcal T}_{\mathbf j}=\sum_{\mathbf i}(-1)^{\sum i_\ell}
       \left(\prod_\ell \mathcal{K}_{j_\ell}^{n_\ell}(i_\ell)\right)A^{\rm split}_{\mathbf i}.
\]
Both transforms are nonnegative. Centralizer inclusion and distance give
\[
 \mathcal T_{\mathbf j}\geq BA^{\rm split}_{\mathbf j},\qquad
 \mathcal T_{\mathbf j}=BA^{\rm split}_{\mathbf j}\quad(1\leq\textstyle\sum j_\ell<d).
\]
These statements hold without assuming purity.

\subsection{Coset-pattern columns}

If a specified isotropic subgroup \(R_i\) lies on block \(i\), enumerate
every coset of \(R_i\) in \(R_i^\perp\). Its pattern
\(p=(p_0,\ldots,p_{n_i})\) counts words at each physical weight.
Let \(y\geq0\) count stabilizer cosets aggregated by their block
patterns and suffix weight. Then a split count is the corresponding
linear combination of products of the \(p_j\)'s.
With no prescribed subgroup a block may instead use the split counts
themselves as variables. Every actual stabilizer containing the prescribed
subgroups gives such nonnegative variables, since it is closed under their
addition. The model need not reconstruct a linear code from an arbitrary
feasible point; only necessity is used.

The row families are:
\begin{center}\small
\begin{tabularx}{\linewidth}{@{}lX@{}}\toprule
Row&Definition\\\midrule
Unit&\(A^{\rm split}_{\mathbf0}=1\).\\
Total&\(\sum A^{\rm split}_{\mathbf i}=B\).\\
Low distance&\(\mathcal T_{\mathbf j}-BA^{\rm split}_{\mathbf j}=0\), \(1\leq\sum j_\ell<d\).\\
Inclusion&\(BA^{\rm split}_{\mathbf j}-\mathcal T_{\mathbf j}\leq0\).\\
Centralizer&\(-\mathcal T_{\mathbf j}\leq0\).\\
Shadow&\(-\widehat{\mathcal T}_{\mathbf j}\leq0\).\\
Weight assumptions&\(\sum_{\sum i_\ell=w}A^{\rm split}_{\mathbf i}=A_w\), or its explicitly assumed lower bound.\\
Coset relations&A specified linear relation valid separately on every physical pattern.\\\bottomrule
\end{tabularx}
\end{center}
The checker verifies each coset relation against the regenerated patterns.
No property of a later shortened code is inserted into these reductions.
The variables may be allowed to take nonnegative real values. This
enlarges the set of possible enumerators, so infeasibility of the
resulting linear system still excludes every actual code. Integrality
is used separately in the explicitly stated parity and counting steps.

\subsection{The sign convention}

Write the system as \(Ez=b\), \(Gz\leq g\), \(z\geq0\).
A certificate consists of arbitrary integer equality multipliers \(\nu\)
and nonnegative integer inequality multipliers \(\lambda\) with
\[
 c=\nu E+\lambda G\geq0,\qquad \nu b+\lambda g<0.
\]
The condition \(c\geq0\) is checked coordinate by coordinate: since
all variables in \(z\) are nonnegative, it makes \(cz\geq0\).
Equality rows may be multiplied by either sign; inequality rows need
nonnegative multipliers so their directions are preserved.
For a feasible \(z\), one would have
\(0\leq cz\leq\nu b+\lambda g<0\), a contradiction.
The checker tests every coefficient column and the negative right-hand side
with exact integers. The multipliers may be found by numerical optimization; the certificate
is checked using integer arithmetic alone.

\subsection{The seventeen-qubit ordinary identity}\label{app:ordinary17}

For the bound on \(A_6+A_7\) in Proposition~\ref{prop:17-enumerator},
write \(\mathcal{K}_j(w)=\mathcal{K}_j^{17}(w)\) and let \(B=2^{14}\). The low-weight rows are
\(\sum_w \mathcal{K}_j(w)A_w-BA_j=0\) for \(1\leq j\leq5\), and let the
shadow rows be \(-\sum_w(-1)^w\mathcal{K}_j(w)A_w\leq0\).
Combine the unit row \(A_0=1\) with multiplier \(-1,835,008\),
the total row with multiplier 116, the five low-weight rows with
\((60,25,11,3,1)\), the five pure-zero rows with
\((-851968,-278528,-507904,-147456,-180224)\), and the shadow rows
\(j=0,\ldots,4\) with \((88,32,18,4,2)\).
The result is
\(32768(A_6+A_7)\leq65536\).
Every coefficient identity is checked independently in the repository.

\subsection{Main certificate inventory}

The repository groups certificates by target parameters. ``Columns'' means
the number of nonnegative variables in the expanded system.

\begin{center}\small
\begin{tabular}{@{}llrr@{}}\toprule
Code&Case&Columns&Negative right-hand side\\\midrule
\(\qc{16,5,5}\)&Weight 1&16&\(-11,520\)\\
&Weight 2&45&\(-374,784\)\\
&Weight 3&42&\(-5,662,720\)\\
&Weight 4, \(A_4\geq2\)&78&\(-7,294,976\)\\
&Weight 5, \(A_1=\cdots=A_4=0\)&72&\(-8,192\)\\
&Minimum stabilizer weight at least 10&17&\(-61,440\)\\\midrule
\(\qc{17,3,6}\)&Weight 1&17&\(-786,432\)\\
&Weight 2&48&\(-5,324,800\)\\
&Weight 3&45&\(-442,368\)\\
&Weight 4&84&\(-81,920\)\\
&Weight 5&78&\(-114,688\)\\
&Pure, \(A_8=0\)&18&\(-3,932,160\)\\\bottomrule
\end{tabular}
\end{center}

The extra seventeen-qubit certificate excluding minimum stabilizer weight
at least nine is also replayed, but is not needed for the forced-weight-eight
route once the \(A_8=0\) certificate has been used. The thirteen simple
local weight-count certificates express the local plane's necessary
weight-six or weight-seven counts; their inequalities are already visible
in Section~\ref{sec:seventeen}.

The fourteen-qubit bundle contains the overlap and rank-three certificates
of Proposition~\ref{prop:14-local}, with its original standalone checker.
The ordinary fourteen-qubit identities and the sixteen-qubit parity
congruence are also replayed separately.
